\section{实验设置}
\subsection{实验任务与候选池}
\subsubsection{固定候选池数据集}
本文采用 \emph{Large Lemma Miners} 公开的 RTL/SVA 基准任务及其可复现工件\cite{peled2026large,peled2026artifact}。每个任务由一个 SystemVerilog 设计和一条待证安全属性构成。沿用该基准的任务划分，主实验组包含 78 个任务，困难实验组包含 31 个任务，共 109 个任务。原工作将至少一种对比模型检查器未能在 10 分钟内证明目标属性的任务归为需要非平凡引理的困难案例\cite{peled2026large}；其余任务构成主实验组。该划分在本文的方法比较前确定，不依据 PS、CH 或 LLM-GLS 的运行结果调整。

为隔离引理选择策略的影响，本文在正式比较前为每个任务生成并冻结一个候选引理池。两个实验组使用相同的非 Agentic 生成配置：模型为 \texttt{gpt-6.1-sol}，推理强度为 \texttt{xhigh}，提供一个 few-shot 示例并采样五次。生成的候选按文本完全一致去重，保留首次出现的顺序并分配候选 ID；任务、属性和逐任务候选池均记录哈希以供核对。PS、CH 和 LLM-GLS 对同一任务读取完全相同的冻结候选池，选择阶段不再生成新候选。

\subsubsection{扩展候选池数据集}
\mbox{}\par
\subsection{对比基线}
\mbox{}\par
\subsection{评价指标与统计方法}
三种方法共用预先生成并冻结的候选池。候选生成耗时不参与选择方法的证明时间比较；证明时间包含选择阶段的 LLM 推理、形式化查询及其他处理开销。结果分析分别报告模型与形式化后端的实际耗时，以检验新增的模型开销是否被节省的验证开销抵消。
\subsection{公平性设置与参数配置}
三种方法采用统一的工具错误终止策略。任一形式化查询因编译错误、后端不支持、进程崩溃或启动失败返回 \texttt{ERROR} 时，当前任务立即结束，并保存错误阶段、候选组合、退出码及后端日志；批次继续处理下一任务。该规则覆盖预检查、候选分类、组合验证、见证提取和 Houdini 搜索，工具错误不触发额外的模型修复。有效反例、\texttt{TIMEOUT} 与 \texttt{UNKNOWN} 分别沿用反例修复及预算管理规则，不与工具错误混同。
